\section{Capitolo 3: Concetti Matematici e Algoritmi}

\begin{esercizio}{Es. 9: Funzione Iniettiva (Verifica Collisioni)}
\textbf{TRACCIA:} Data una funzione `f(x)` e un dominio (array di input) `D`, scrivere `isIniettiva(f, D)` che restituisce `true` se la funzione non produce mai lo stesso output per due input diversi in D.
\end{esercizio}

\begin{jsbox}
function isIniettiva(f, dominio) {
    let immaginiViste = new Set();
    
    for (let x of dominio) {
        let y = f(x); // Calcolo immagine
        
        if (immaginiViste.has(y)) {
            return false; // Collisione trovata! (Due x danno la stessa y)
        }
        immaginiViste.add(y);
    }
    return true; // Nessuna collisione
}
\end{jsbox}

\begin{esercizio}{Es. 10: Convoluzione Discreta (Semplificata)}
\textbf{TRACCIA:} Dati due array `A` e `B` della stessa lunghezza, calcolare il prodotto scalare (somma dei prodotti elemento per elemento).
$S = \sum (A[i] \cdot B[i])$
\end{esercizio}

\begin{jsbox}
function prodottoScalare(A, B) {
    let somma = 0;
    // Assumiamo A.length === B.length
    for (let i = 0; i < A.length; i++) {
        somma += A[i] * B[i];
    }
    return somma;
}
\end{jsbox}

\begin{esercizio}{Es. 11: Punto Fisso}
\textbf{TRACCIA:} Scrivere `trovaPuntoFisso(f, min, max)` che cerca se esiste un numero intero $x$ tra `min` e `max` tale che $f(x) == x$. Se esiste ritornalo, altrimenti `null`.
\end{esercizio}

\begin{jsbox}
function trovaPuntoFisso(f, min, max) {
    for (let x = min; x <= max; x++) {
        if (f(x) === x) {
            return x; // Trovato!
        }
    }
    return null;
}
\end{jsbox}

\begin{esercizio}{Es. 12: Numeri Triangolari Ricorsivi}
\textbf{TRACCIA:} L'n-esimo numero triangolare è la somma dei numeri da 1 a n. Scrivere una funzione ricorsiva.
\end{esercizio}

\begin{jsbox}
function triangolare(n) {
    if (n <= 0) return 0;
    return n + triangolare(n - 1);
}
// Esempio: T(3) = 3 + T(2) = 3 + 2 + T(1) = 3+2+1 = 6
\begin{esercizio}{Es. 23: Numero Primo (IsPrime)}
\textbf{TRACCIA:} Scrivere `isPrime(n)` che torna true se $n$ è primo.
\textit{Ricorda: 1 non è primo. Basta controllare fino a $\sqrt{n}$.}
\end{esercizio}

\begin{jsbox}
function isPrime(n) {
    if (n < 2) return false;
    
    // Ottimizzazione: fino alla radice quadrata
    for (let i = 2; i * i <= n; i++) {
        if (n % i === 0) return false; // Divisibile
    }
    return true;
}
\end{jsbox}

\begin{esercizio}{Es. 24: MCD (Algoritmo di Euclide)}
\textbf{TRACCIA:} Calcolare il Massimo Comune Divisore tra `a` e `b`.
\textit{Algoritmo: Se b=0, MCD è a. Altrimenti MCD(b, a \% b).}
\end{esercizio}

\begin{jsbox}
function mcd(a, b) {
    if (b === 0) return a;
    return mcd(b, a % b);
}
\end{jsbox}
